\documentclass[conference]{IEEEtran}

\usepackage{amsmath,graphicx}
\usepackage{url} 

\title{Capstone Report - Steam Player Retention Project \\[0.2em]
\large Predictive Modeling for Player Retention in Multi-Modal Gaming Data}

\author{
\IEEEauthorblockN{Jason Limfueco, Laura Ngo, Baixi Guo}
\IEEEauthorblockA{
University of California San Diego, Master of Data Science
}
}

\begin{document}

\maketitle

\begin{abstract}
Steam is one of the largest digital gaming platforms in the world, generating vast amounts of multi-modal data that capture player behavior, game interactions, and community sentiment at scale. Player retention is a critical factor in the success of these online gaming platforms, as it directly influences user engagement, personalization, and long-term sustainability and revenue. In this project, we define player retention as a supervised binary classification problem, where the goal is to predict whether a user will remain active on the platform given historical game play behavior, user interaction patterns, and game metadata. The input consists of structured behavioral logs, textual data (e.g., reviews), and derived semantic features, while the output is a binary retention label. We leverage a large-scale Steam dataset (43GB), reducing the high-dimensional feature space through embedding techniques to support both predictive modeling and exploratory analysis \cite{steam_dataset}. Principal Component Analysis (PCA) is applied to compress feature representations, balancing information retention with computational efficiency. Textual features are represented using TF-IDF embeddings, which are paired with nearest neighbor methods to support fast and scalable recommendation generation. 

Prior work in retention modeling has primarily focused on behavioral features in domains such as e-commerce and gaming. For example, recent studies have applied traditional machine learning models to play churn prediction to exemplify the effectiveness of game-play and interaction data, while highlighting limitations in data diversity \cite{peng2024churn}. More advanced approaches explored deep learning methods and improved predicting performance through capturing complex temporal and behavioral patterns in player activity \cite{hoang2025}. However, these approaches lack integration of heterogeneous data sources such as textual user feedback. Our approach extends this line of work by incorporating sentiment analysis and embedding-based representations to capture richer user and game dynamics. 

We evaluate a range of models, including Logistic Regression as an interpretable baseline, and more expressive methods such as Random Forests and Gradient Boosted Trees to capture nonlinear relationships. Additionally, we explore temporal patterns and embedding-based features to further improve predictive performance. Beyond predictive modeling, we adopt a multi-stage analytical framework. Our system involves predicting player retention using supervised learning methods, analyzing misclassifications to identify factors contributing to model error and uncover nuanced user behavior patterns, and utilizing a similarity-based recommendation system to identify competing games and potential alternatives influencing player churn. 

Model performance is assessed using standard classification metrics including accuracy, precision, recall, F1-score, and ROC-AUC, to evaluate predictive performance and model robustness. Beyond predictive accuracy, we aim to provide actionable insights into the key factors that drive player retention in gaming ecosystems. This paper presents the analytical framework, system architecture, and visualizations developed to address the challenge of modeling and understanding player retention and engagement on Steam. 
\end{abstract}

\begin{thebibliography}{1}

\bibitem{steam_dataset}
K. Pocock, "Steam Reviews Dataset," Kaggle, 2023. [Online]. Available: \url{https://www.kaggle.com/datasets/kieranpoc/steam-reviews/code}. [Accessed: Apr. 8, 2026].

\bibitem{peng2024churn}
X. Peng et al., "Research on Machine Learning Models for Predicting Player Churn," Proc. Int. Conf. Data Science and Applications, 2024.

\bibitem{hoang2025}
T. Hoang et al., "Early-Stage Churn Prediction in Online Games Using Neural Networks," Engineering Applications of Artificial Intelligence, 2025.

\end{thebibliography}

\end{document}
